%%% FIELD claim-finite-stratum-compiler
For each of the at most \(N!\) topological orders on \(W\), one saturated superarchitecture containing every forward directed slot and one potential source \(U_C\) for every nonempty \(C\subseteq W\) contains every finite retained-variable response architecture having that order: unused directed inputs are made inessential, and an absent source has a canonical embedding as a simplex-vertex source ignored by every incident table.  Sources with the same realized nonempty child set can be merged without changing any structural solution, so at most \(H_{\max}=2^N-1\) active independent sources are needed.  On any fixed permitted architecture, sourcewise barycentric reduction preserves \(\zeta(M)\) with at most \(K_{\max}=m+3\) values per source, after which the model embeds in its saturated ordered superarchitecture and its dense deterministic tables are covered by \(R_{\max}\).  Every point of \(\mathfrak F_\delta(\mathcal P,q)\) induces an SCM in \(\mathfrak M_\delta(\mathcal P,q)\), and the target-value image of \(\mathfrak F_\delta(\mathcal P,q)\) is exactly the intervention-value set of the enumerated bounded finite SCMs satisfying actual selected-MAG legality.  Sourcewise barycentric reduction is asserted to preserve moments, not strict legality; no assertion identifies the bounded image with the image of all arbitrary admissible SCMs.
%%% FIELD claim-certified-attainment-procedure
For every finite rational instance, the certified procedure uses at most \(A_{\max}=N!\) saturated ordered superarchitectures and, for each order, at most \(R_{\max}\) dense deterministic response-table systems.  It expands the moment and essential-feature polynomials coefficient by coefficient and reports INFEASIBLE or isolates the algebraic infimum and supremum of the bounded finite graph-legal image.  At either isolated closure endpoint \(t\), it decides \(\mathsf{Lift}(t)\).  If true, \(\mathsf{Cert}(t)\) contains a branch \((T,G)\) and a real-univariate representation with Thom encoding of algebraic \(p_T,\rho\) whose exact graph is verified by \(\Psi_{T,G}\).  If false, \(\mathsf{Cert}(t)\) contains, for every enumerated bounded branch, real-univariate sample points and their complete sign lists meeting every semialgebraically connected component of every realizable sign condition and proving that the branch is empty; a sign-invariant cylindrical decomposition is an optional alternative certificate.  For one existential block of \(p_{\mathrm{dim}}\) response parameters and one free endpoint variable, with \(s_{\mathrm{count}}\) input polynomials, maximum degree \(D_{\mathrm{deg}}\), and expanded integer coefficient bit length \(\tau_{\mathrm{eff}}\), block elimination uses at most
\[
A_{\max}R_{\max}s_{\mathrm{count}}^{\,2(p_{\mathrm{dim}}+1)}D_{\mathrm{deg}}^{O(p_{\mathrm{dim}})}
\]
ring operations, produces univariate output degree \(D_{\mathrm{deg}}^{O(p_{\mathrm{dim}})}\), and has intermediate and output integer bit lengths \(\tau_{\mathrm{eff}}D_{\mathrm{deg}}^{O(p_{\mathrm{dim}})}\).  In coarse bit-operation notation the bound is
\[
A_{\max}R_{\max}(s_{\mathrm{count}}D_{\mathrm{deg}})^{O(p_{\mathrm{dim}})}\operatorname{poly}(\tau_{\mathrm{eff}}).
\]
A negative endpoint certificate excludes only the enumerated bounded finite compiler branches and does not resolve universal legal lifting for arbitrary admissible SCMs.
%%% FIELD construction-architecture-list
For each topological order on \(W\), form one saturated superarchitecture containing every forward directed slot and one formal independent source \(U_C\) for every nonempty \(C\subseteq W\).  Thus at most \(A_{\max}=N!\) labeled superarchitectures are generated.  Every source has at most \(K_{\max}=m+3\) slots, including zero-mass slots, and every dense deterministic response table is enumerated, with at most \(R_{\max}\) table systems per order.  To embed a sparse architecture without loss, make each added directed input inessential.  Represent an absent source by a simplex vertex, with one slot of mass one and all remaining slots of mass zero, and make every incident table ignore it; an all-zero probability vector is never used.  This simplex-vertex convention is only a canonical embedding of a missing source and is not imposed on every parameter point having no active source incidence.  If a formally named \(U_C\) is active only on \(D\subsetneq C\), its realized child set is \(D\).  Before graph projection, discard empty realized child sets and merge formal sources with the same nonempty realized child set into one tuple source.  Retain precisely the exact projected selected MAGs \(G\in\mathcal G(\mathcal P)\).
%%% FIELD construction-graph-legality
For a saturated order \(T\), fixed deterministic tables, and masses \(p_T\), retain every directed sign \(Z_{\operatorname{dir}(i,j)}(T,p_T)\) and every formal source-incidence sign \(Z_{\operatorname{src}(C,j)}(T,p_T)\).  A directed slot is actual exactly when its sign is positive.  For each formal \(U_C\), define solely from its active incidence signs
\[
D_C(p_T)=\{j\in C:Z_{\operatorname{src}(C,j)}(T,p_T)>0\}.
\]
Discard a source with \(D_C(p_T)=\varnothing\) when forming the actual graph, without constraining its mass vector to a simplex vertex and without requiring its literal table entries to be constant: the zero incidence signs already say that any variation is functionally inessential outside exogenous null sets.  Replace every remaining formal label \(C\) by its realized nonempty child set \(D_C(p_T)\), and merge sources with equal realized child sets into one tuple source.  Apply the finite latent-and-selection projection rules to the resulting actual directed edges and normalized nonempty sources.  The formula \(\Psi_{T,G}(p_T)\) is the quantifier-free Boolean combination of the complete positive-or-zero incidence sign pattern and the finite projection test; it is true exactly when \(G_S(M_{T,p_T})=G\).  Every endogenous table comparison in \(Z_e\) ranges over all table inputs, including off-support endogenous assignments.
%%% FIELD construction-certified-procedure
Generate the at most \(A_{\max}=N!\) saturated ordered superarchitectures, enumerate their dense deterministic response-table systems up to \(R_{\max}\), and expand every coordinate of \(\Phi_T\) and every \(Z_e\) coefficient by coefficient as a rational polynomial.  Introduce \(t\) through \(\Phi_T(p_T)_n-t\Phi_T(p_T)_d=0\), and form \(\mathsf{Lift}(t)\) as the finite disjunction over table systems and \(G\in\mathcal G(\mathcal P)\).  Use the one-block existential quantifier-elimination algorithm in \(\text{lem:basu-one-block-elimination}\) to report INFEASIBLE or to produce a univariate sign formula and isolate the algebraic infimum and supremum of the bounded finite image.  At an isolated closure endpoint, decide \(\mathsf{Lift}(t)\) using \(\text{lem:basu-sign-sampling}\).  The certificate \(\mathsf{Cert}(t)\) is either a real-univariate representation with Thom encoding of a satisfying legal response point, or, for every bounded table branch, a finite set of real-univariate sample points with complete sign lists meeting every semialgebraically connected component of every realizable sign condition and proving the branch empty.  A sign-invariant cylindrical decomposition may instead be retained as an optional certificate format, but it is not used for the intrinsic complexity bound.
%%% FIELD lemma-basu-sign-statement
Let \(Q\) be a polynomial in \(p_{\mathrm{dim}}\) variables, let \(\mathscr P\) contain at most \(s_{\mathrm{count}}\) polynomials in the same variables, suppose that all these polynomials have degree at most \(D_{\mathrm{deg}}\), and let \(p_0\) denote the real dimension of the zero set of \(Q\).  Over a real closed coefficient field, there is an algorithm computing a finite set meeting every semialgebraically connected component of every realizable sign condition of \(\mathscr P\) over that zero set in \(s_{\mathrm{count}}^{p_0}D_{\mathrm{deg}}^{O(p_{\mathrm{dim}})}\) ring operations.  Returning the signs of every polynomial at every sample costs \(s_{\mathrm{count}}^{p_0+1}D_{\mathrm{deg}}^{O(p_{\mathrm{dim}})}\) ring operations.  The algebraic samples may be recorded by real-univariate representations whose selected roots are specified by Thom encodings.  Taking \(Q=0\) and \(p_0=p_{\mathrm{dim}}\) gives the ambient-space bounds used here.
%%% FIELD lemma-basu-block-statement
Let a first-order formula over a real closed field have at most \(s_{\mathrm{count}}\) input polynomials of degree at most \(D_{\mathrm{deg}}\), one block of \(p_{\mathrm{dim}}\) quantified variables, and one free variable.  It has an equivalent quantifier-free formula computable in
\[
s_{\mathrm{count}}^{\,2(p_{\mathrm{dim}}+1)}D_{\mathrm{deg}}^{O(p_{\mathrm{dim}})}
\]
ring operations.  Its output polynomials have degree \(D_{\mathrm{deg}}^{O(p_{\mathrm{dim}})}\).  For integer input coefficients of bit length at most \(\tau_{\mathrm{eff}}\), all intermediate and output integer bit lengths are at most \(\tau_{\mathrm{eff}}D_{\mathrm{deg}}^{O(p_{\mathrm{dim}})}\).
%%% FIELD proof-finite-stratum-compiler
We prove the assertions in four steps.

First fix a topological order of \(W\).  By \(\text{def:architecture-list}\), its saturated superarchitecture contains every forward directed argument and every formal source \(U_C\), \(\varnothing\ne C\subseteq W\).  Consider any finite retained-variable response architecture having this order.  Replace each missing directed argument by a table argument on which the corresponding structural function is constant.  If a source is absent, put its mass vector at a simplex vertex, with one slot of mass one and all other slots of mass zero, and make every incident table constant in that source argument.  Thus no all-zero mass vector is used.  If a source has actual child set \(D\), install its law and table arguments in the formal slot \(U_D\), after applying the merging operation proved next if necessary.  These extensions change no effective structural equation and hence change no factual or interventional solution.  Every formerly sparse finite architecture therefore embeds in the saturated architecture for its order.  There are at most \(N!\) orders.

We next justify source merging and the realized-child-set convention.  Group mutually independent exogenous sources according to their nonempty sets of actual retained children.  For each \(\varnothing\ne C\subseteq W\), let \(U_C\) be the tuple of all original sources with actual retained child set \(C\), equipped with the product law inherited from their mutual independence.  Tuples for distinct \(C\)'s use disjoint original coordinates and hence remain mutually independent.  Substitute the tuple coordinates into the old structural equations.  Evaluation in topological order proves, node by node, that every structural solution under every intervention is unchanged.  A retained equation depends essentially on the tuple precisely when it depends essentially on at least one tuple coordinate.  One direction is immediate.  For the converse, join two tuple values by changing one coordinate at a time; if the endpoint outputs differ, the output changes at some step.  Thus tuple merging preserves the minimal structural graph.  Sources with no actual retained child may be discarded.  There are exactly \(2^N-1=H_{\max}\) nonempty subsets of \(W\), establishing the active-source bound.

For an arbitrary parameter point of a saturated superarchitecture, define
\[
D_C(p_T)=\{j\in C:Z_{\operatorname{src}(C,j)}(T,p_T)>0\}.
\]
By \(\text{lem:essential-feature-iff}\), \(D_C(p_T)\) is exactly the actual retained child set of the formal source \(U_C\).  Consequently a formal source with empty \(D_C(p_T)\) is absent from the minimal structural graph and is discarded for projection, even if its mass vector is nondegenerate or its literal table varies only on null exogenous contexts.  No normalization constraint is imposed on that parameter point.  When \(D_C(p_T)\ne\varnothing\), relabel the source by \(D_C(p_T)\); tuple-merging all formal sources with the same realized nonempty child set preserves solutions and the minimal graph by the preceding paragraph.  Therefore the corrected formula \(\Psi_{T,G}(p_T)\), which projects only after empty-set deletion and equal-set merging, is equivalent to \(G_S(M_{T,p_T})=G\).  This is an interpretation of the active incidence signs, not an extra restriction on source masses or tables.

Second, fix any permitted response architecture with independent sources \(U_1,\ldots,U_H\), where \(H\leq H_{\max}\).  Define
\[
F(U_1,\ldots,U_H)=
\left(
  (\mathbf 1\{V=v,S=1\})_{v\in\Omega_V},
  h(A^b)\mathbf 1\{S^b=1\},
  \mathbf 1\{S^b=1\}
\right)\in\mathbb R^{m+2}.
\]
The boundedness of \(h\) and finiteness of \(\Omega_V\) make \(F\) integrable, and \(\mathbb E_M F=\zeta(M)\) by \(\text{def:moment-map}\).  Hold all source laws except that of \(U_k\) fixed and define
\[
g_k(u)=
\int F(u_1,\ldots,u_{k-1},u,u_{k+1},\ldots,u_H)
      \prod_{r\ne k}dP_r(u_r).
\]
Mutual independence and iterated integration give
\[
\zeta(M)=\int g_k(u)\,dP_k(u).
\]

We record the elementary barycenter argument used here.  If an integrable random vector \(Y\) takes values in \(\mathbb R^d\), then \(\mathbb EY\) belongs to the convex hull of its essential range.  Approximation by integrable simple random vectors first places \(\mathbb EY\) in the closure of that convex hull.  If it were not in the convex hull itself, then it would lie on the relative boundary of the closed convex hull: a point in the relative interior of a convex set's closure is in the relative interior of the convex set.  A supporting affine functional at that boundary is nonnegative relative to its value at \(\mathbb EY\).  Its nonnegative difference has expectation zero and hence vanishes almost surely.  Restricting the essential range to the resulting proper affine hyperplane and inducting on affine dimension proves convex-hull membership.  By the definition of convex hull there is then a finite convex representation.  If it uses more than \(d+1\) points, an affine dependence among them permits changing the weights, while preserving their sum and barycenter, until one weight vanishes.  Iteration leaves at most \(d+1\) actual essential-range values.

Apply this argument with \(d=m+2\).  There are \(L_k\leq m+3=K_{\max}\), values \(u_{k,1},\ldots,u_{k,L_k}\), and weights \(w_{k,r}\geq0\) such that
\[
\sum_{r=1}^{L_k}w_{k,r}=1,
\qquad
\sum_{r=1}^{L_k}w_{k,r}g_k(u_{k,r})
=\int g_k(u)\,dP_k(u).
\]
Replacing only \(P_k\) by this discrete law preserves \(\zeta(M)\) and product factorization.  Repeat for \(k=1,\ldots,H\), always using the currently installed laws of the other sources.  Every replacement preserves the whole moment vector, so the terminal model still has moment vector \(\zeta(M)\) and every source has at most \(K_{\max}\) values.  In particular,
\[
\rho=\sum_{v\in\Omega_V}c_v(M)
\]
is preserved as one common normalization.  Since retained state spaces and reduced source supports are finite, every structural function is a deterministic response table.  Embedding the reduced architecture in its saturated order only adds inessential directed inputs and canonical simplex-vertex ignored sources, so its dense table system is among those counted by \(R_{\max}\).  This reduction does not assert that any strict inequality \(Z_e>0\) survives.

Third, take \((p_T,\rho)\in\mathfrak F_\delta(\mathcal P,q)\).  Its saturated topological order, deterministic tables, and mutually independent source simplexes define a recursive finite SCM \(M_{T,p_T}\).  The graph-legality stratum determines active incidence signs, discards empty realized source child sets, and merges duplicates only for interpreting the actual graph.  The essential-feature lemma gives a \(G\in\mathcal G(\mathcal P)\) such that
\[
G_S(M_{T,p_T})=G.
\]
The remaining constraints in \(\text{def:finite-legal-fiber}\) give, for each \(v\in\Omega_V\),
\[
c_v(M_{T,p_T})=\Phi_T(p_T)_{c_v}=\rho q(v),
\qquad
\delta\leq\rho\leq1,
\qquad
d_{M_{T,p_T}}=\Phi_T(p_T)_d\geq\delta.
\]
Together with \(\text{ass:positive-selected-law}\), these are exactly the member conditions in \(\text{def:admissible-class}\).  Hence \(M_{T,p_T}\in\mathfrak M_\delta(\mathcal P,q)\).  Since \(\delta>0\), its denominator is positive and
\[
\theta(M_{T,p_T})=
\frac{\Phi_T(p_T)_n}{\Phi_T(p_T)_d}
\]
is well defined.

Finally, every enumerated bounded finite SCM with actual selected MAG in \(\mathcal G(\mathcal P)\) has a topological order and embeds in the corresponding saturated superarchitecture by the first step.  Its deterministic table system is counted by \(R_{\max}\), its product source masses define \(p_T\), and \(\rho=\sum_vc_v\).  Selected-law fit and the two margins put \((p_T,\rho)\) in \(\mathfrak F_\delta(\mathcal P,q)\), while the essential-feature lemma selects its exact graph branch.  The two inclusions prove equality of the two bounded finite target-value images.  None of these steps proves the legal lifting asked for in \(\text{oeq:universal-legal-lifting}\).
%%% FIELD proof-certified-attainment-procedure
Fix one saturated order and one deterministic response-table system.  Let its \(H\leq H_{\max}\) formal source-mass vectors be \(p_1,\ldots,p_H\).  For every coordinate \(y\) of the moment map there are rational table coefficients \(a_y(u_1,\ldots,u_H)\) for which
\[
\Phi_T(p_T)_y
=\sum_{u_1,\ldots,u_H}
  a_y(u_1,\ldots,u_H)\prod_{k=1}^H p_k(u_k).
\]
For a selected cell or \(d\), each coefficient is zero or one; for \(n\), it is either zero or a value of the rational map \(h\).  Thus every moment coordinate is a rational polynomial of degree at most \(H\).  By \(\text{lem:essential-feature-iff}\), a directed-feature polynomial integrates a zero-one table indicator over at most one copy of every source and has degree at most \(H\).  A source-feature polynomial uses two values of its own source and at most one value of every other source, so its degree is at most \(H+1\).  Simplex normalization and nonnegativity, the equations
\[
\Phi_T(p_T)_{c_v}=\rho q(v)\qquad(v\in\Omega_V),
\]
the inequalities \(\delta\leq\rho\leq1\) and \(\Phi_T(p_T)_d\geq\delta\), the complete active-incidence sign conditions used by \(\Psi_{T,G}\), and
\[
\Phi_T(p_T)_n-t\Phi_T(p_T)_d=0
\]
are consequently a finite Boolean combination of rational polynomial sign conditions.  The realized child sets \(D_C\), deletion of empty sets, merging of equal nonempty sets, and selected projection depend only on that finite complete sign pattern.  They impose no extra equation requiring an inessential formal source to be degenerate or its table to be literally constant.  The declared rationality of \(q\), \(\delta\), and \(h\) permits clearing positive common denominators without changing any sign.  By definition, \(s_{\mathrm{count}}\), \(D_{\mathrm{deg}}\), and \(\tau_{\mathrm{eff}}\) bound, respectively, the number, degree, and expanded integer coefficient length of the resulting polynomials.

The saturated architecture definition gives at most \(A_{\max}=N!\) orders.  For a fixed order, every retained equation has all forward directed arguments and every potentially incident source argument, so the dense deterministic table enumeration is bounded by \(R_{\max}\).  A sparse architecture has a canonical embedding using inessential added directed inputs and simplex-vertex ignored absent sources, but the table branch itself ranges over every enumerated mass vector and table.  Absence and activity for an arbitrary parameter point are read only from the essential-feature signs.  Thus missing arrows and sources need not be enumerated as separate architectures, and there are at most \(A_{\max}R_{\max}\) table branches.

For each branch, all response parameters, including \(\rho\), form one existential block of size at most \(p_{\mathrm{dim}}\), while \(t\) is one free variable.  Apply \(\text{lem:basu-one-block-elimination}\), namely Basu's Theorem 2.27 and its preceding block-elimination discussion, with \(\omega=1\), \(k_1=p_{\mathrm{dim}}\), and \(\ell=1\).  It returns an equivalent quantifier-free formula in \(t\).  Its ring-operation cost is
\[
s_{\mathrm{count}}^{\,2(p_{\mathrm{dim}}+1)}
D_{\mathrm{deg}}^{O(p_{\mathrm{dim}})},
\]
the output degree is \(D_{\mathrm{deg}}^{O(p_{\mathrm{dim}})}\), and integer intermediate and output bit lengths are bounded by
\[
\tau_{\mathrm{eff}}D_{\mathrm{deg}}^{O(p_{\mathrm{dim}})}.
\]
Multiplication by the number of table branches gives the displayed ring-operation bound.  Exact integer arithmetic on operands of the displayed bit lengths, together with exact univariate root isolation, is polynomial in those bit lengths and in the univariate output size; it is therefore absorbed in the coarser bit-operation expression
\[
A_{\max}R_{\max}
(s_{\mathrm{count}}D_{\mathrm{deg}})^{O(p_{\mathrm{dim}})}
\operatorname{poly}(\tau_{\mathrm{eff}}).
\]
Theorem 2.18 of the cited survey is a bound on connected components, not a quantifier-elimination theorem, and is not used.

The strict feature inequalities remain strict in the formula, so its existential projection is the exact bounded branch of \(\mathsf{Lift}(t)\), not the closure of that branch.  If no projected branch is nonempty, the procedure reports INFEASIBLE.  Otherwise every feasible endpoint variable is bounded.  Indeed, pointwise on \(\{S^b=1\}\),
\[
h_-\leq h(A^b)\leq h_+,
\]
and multiplication by \(\mathbf 1\{S^b=1\}\) followed by expectation gives
\[
h_-d\leq n\leq h_+d.
\]
The intervention margin gives \(d\geq\delta>0\); therefore the endpoint equation implies
\[
h_-\leq t=\frac nd\leq h_+.
\]
The projected set is a bounded univariate semialgebraic set.  Between consecutive real roots of its finitely many output polynomials every sign is constant.  Its infimum and supremum must therefore be either \(h_-\), \(h_+\), or a real root of an output polynomial.  Hence they are real algebraic, have degree \(D_{\mathrm{deg}}^{O(p_{\mathrm{dim}})}\), and can be isolated by a Thom encoding.

Fix either isolated endpoint \(t_0\), and regard its Thom encoding as specifying the ordered real algebraic coefficient extension obtained after substitution.  Apply \(\text{lem:basu-sign-sampling}\), Basu's Theorem 2.16 with Definitions 2.12--2.13, to the branch polynomials.  Specializing its algebraic-set polynomial to the zero polynomial makes the ambient zero set the full \(p_{\mathrm{dim}}\)-dimensional response space.  It produces a finite set of real-univariate sample points meeting every semialgebraically connected component of every realizable sign condition.  Asking also for all signs costs
\[
s_{\mathrm{count}}^{\,p_{\mathrm{dim}}+1}
D_{\mathrm{deg}}^{O(p_{\mathrm{dim}})}
\]
ring operations, which is dominated by the preceding one-block elimination bound after enlarging the absolute constants.

If some sample sign vector satisfies all simplex, moment, margin, endpoint, and graph-sign clauses, its real-univariate representation and Thom encoding give algebraic \(p_T,\rho\).  By \(\text{lem:essential-feature-iff}\), the signs in \(\Psi_{T,G}\) verify
\[
G_S(M_{T,p_T})=G\in\mathcal G(\mathcal P),
\]
so this is the asserted positive certificate.

Conversely, suppose no sample sign vector satisfies a given branch.  If a response point satisfied it, that point would lie in a semialgebraically connected component of a realizable sign condition of the branch polynomials.  The sample set meets that component, and the defining Boolean formula depends only on the complete sign vector; its sample would therefore satisfy the same branch, a contradiction.  Hence the branch is empty.  Collecting these sample representations and sign lists for all finitely many branches proves \(\neg\mathsf{Lift}(t_0)\).  A full sign-invariant cylindrical decomposition may be retained instead, but its potentially doubly exponential construction is not the intrinsic complexity asserted here.  Finally, every branch quantified in this argument belongs to the bounded finite compiler.  Thus a negative certificate excludes all and only those enumerated bounded branches; without an affirmative resolution of \(\text{oeq:universal-legal-lifting}\), it does not exclude an equal-moment model in the arbitrary-SCM superclass.
%%% FIELD obligation-partial-result
The mandatory collider test closes for the following nontrivial subcase.  Let \(V=\{X,Y,Z\}\) be binary, let the retained structural graph be \(X\to Z\leftarrow Y\) with \(S\) an independent selection root, take \(B=\{X\}\), and assume that the positive selected law \(q\) obeys the collider Markov restriction \(X\mathbin{\perp\!\!\!\perp}Y\).  Put
\[
p_{xy}=q(Z=1\mid X=x,Y=y)\in(0,1).
\]
Choose independent \(U_X,U_Y\) with laws \(q_X,q_Y\), an independent \(U_Z\sim\operatorname{Unif}[0,1]\), and measurable sets \(E_{xy}\subset[0,1]\) with Lebesgue measure \(p_{xy}\).  The four sets can be chosen so that \(E_{0y}\ne E_{1y}\) modulo null sets for at least one \(y\), and \(E_{x0}\ne E_{x1}\) modulo null sets for at least one \(x\): begin with intervals of the prescribed lengths, and whenever an equality remains, exchange two sufficiently short subintervals between the set and its complement.  Positivity \(0<p_{xy}<1\) makes each exchange possible, and finitely many disjoint exchanges preserve distinctions already installed.  Define
\[
X=U_X,\qquad Y=U_Y,\qquad Z=\mathbf 1\{U_Z\in E_{XY}\}.
\]
Then
\[
\Pr(X=x,Y=y,Z=z)=q_X(x)q_Y(y)q(z\mid x,y)=q(x,y,z),
\]
and the two strict symmetric differences make both arrows into \(Z\) actual functional dependencies, even when \(q\) is unfaithful to one or both arrows.  Taking an independent \(S\sim\operatorname{Bernoulli}(\rho)\), with \(\rho\in[\delta,1]\), gives \(c_v=\rho q(v)\) and \(d=\rho\).  For every bounded binary outcome map \(h\), truncated factorization gives the single attainable selected-intervention value
\[
\theta=
\sum_{y\in\{0,1\}}q_Y(y)
\bigl[h(0)(1-p_{by})+h(1)p_{by}\bigr].
\]
It is attained by the exact collider construction above.  Thus the exact-graph fiber and all its support-stratum closure points lift in this isolated-selection compelled-collider case.  This does not cover selected MAGs in which \(S\) has parents or latent ancestors, nor does it preserve an arbitrary pair \((n,d)\) beyond the value fixed by this collider factorization.
